\chapter{The Growth of Functions: Asymptotic Notation}

In the last chapter you counted the work that insertion sort does. You found
that, in the worst case, the number of basic steps came out to something like
$\frac{1}{2}n^2 + \frac{5}{2}n - 3$. That is an exact count. And yet,
if you carry it around forever, it will weigh you down. The constant $\frac12$
depends on how you decided to count comparisons. The $\frac52 n$ term depends on
bookkeeping you might do differently tomorrow. The $-3$ is an accident of where
the loops start. None of those details survive a change of computer, compiler,
or mood.

What \emph{does} survive is the shape of the function: it grows like $n^2$. Double
the input and the running time roughly quadruples. \emph{That} is the durable
truth, and this chapter is about the language we use to state it precisely. The
language is called \textbf{asymptotic notation}, and once it clicks you will read
running times the way a musician reads a key signature: at a glance, without
sounding out every note.

\begin{focusbox}
The single idea of this chapter: when comparing algorithms, we care about
\emph{how the running time grows as the input grows}, not about its exact value
on small inputs. Asymptotic notation ($O$, $\Omega$, $\Theta$, and their little
cousins) is the vocabulary for talking about growth rate while deliberately
ignoring constant factors and low-order terms.
\end{focusbox}

\section{Why we throw away the constants}

Suppose Algorithm A runs in exactly $100n$ steps and Algorithm B runs in exactly
$2n^2$ steps. Which is faster? The answer is: \emph{it depends on $n$}.

Let us actually compute. For tiny inputs, B wins because A pays a large fixed toll:

\begin{center}
\begin{tabular}{rrr}
\toprule
$n$ & $100n$ (A) & $2n^2$ (B) \\
\midrule
1    & 100      & 2 \\
10   & 1{,}000  & 200 \\
50   & 5{,}000  & 5{,}000 \\
100  & 10{,}000 & 20{,}000 \\
1000 & 100{,}000 & 2{,}000{,}000 \\
\bottomrule
\end{tabular}
\end{center}

Up to $n=50$ they are neck and neck, and below that B is the better deal. But look
what happens after the crossover. At $n=1000$, A does $100{,}000$ steps while B does
two million. Push to $n=10^6$ and A does $10^8$ steps while B does $2\times 10^{12}$
--- B is now twenty thousand times slower. The gap does nothing but widen, forever.

Figure~\ref{ch03:fig-crossover} shows the shape of that story. The quadratic B starts
below the linear A, but the two curves cross once and only once; past the crossover
$n^\ast$, B is above A and pulls away with no hope of return.

\begin{figure}[H]
\centering
\begin{tikzpicture}[scale=1.0]
  \draw[->] (0,0) -- (6.4,0) node[right] {$n$};
  \draw[->] (0,0) -- (0,4.4) node[above] {steps};
  % A: linear 100n, drawn gently
  \draw[thick,sourceblue,domain=0:6,smooth,variable=\x]
        plot ({\x},{0.55*\x}) node[right]{$100n$ (A)};
  % B: quadratic 2n^2, starts low, crosses, shoots up
  \draw[thick,trapred,domain=0:2.85,smooth,variable=\x]
        plot ({\x},{0.40*\x*\x}) node[above]{$2n^2$ (B)};
  % crossover point: 0.55 x = 0.40 x^2 -> x = 1.375
  \draw[dashed,gray] (1.375,0) -- (1.375,0.756);
  \fill[black] (1.375,0.756) circle (1.4pt);
  \node[below] at (1.375,0) {$n^\ast$};
  \node[gray] at (4.5,0.55) {\small before $n^\ast$: B wins};
  \node[gray] at (2.7,3.4) {\small after $n^\ast$: A wins forever};
\end{tikzpicture}
\caption{Linear A versus quadratic B. A pays a large constant toll and so loses for small
$n$, but a single crossover at $n^\ast$ flips the order permanently: the faster-growing
curve eventually dominates regardless of the constants.}
\label{ch03:fig-crossover}
\end{figure}

\begin{intuitionbox}
The constant in front ($100$ versus $2$) buys you an advantage on small inputs,
but it is a \emph{one-time} advantage that gets paid off. The exponent --- linear
versus quadratic --- compounds. Asymptotic analysis is the decision to bet on the
exponent, because for any fixed constant there is some input size past which the
slower-growing function wins and never looks back. We care about ``eventually and
forever,'' not ``on my laptop today with $n=12$.''
\end{intuitionbox}

There is a second reason to drop constants, and it is humbler: \emph{we usually do
not know them}. The exact constant depends on the machine's clock, the cache, the
compiler's cleverness, whether the data sits in registers. If we promised a precise
step count we would be lying about precision we do not have. By contrast, the growth
rate --- $n$ versus $n^2$ versus $\log n$ --- is a property of the \emph{algorithm},
not the machine. It is the part worth naming.

\begin{aestheticbox}
This is not academic hair-splitting; it is how working engineers triage. When a
service slows to a crawl as data grows, the first question is never ``what is the
constant'' but ``what is the growth rate.'' An $O(n^2)$ routine that was invisible
at a thousand records will melt the server at a million; an $O(n\log n)$ routine
will shrug. Choosing the right growth class is the difference between a feature
that scales and one that has to be rewritten at 2 a.m. The notation in this chapter
is the shorthand every performance discussion is conducted in.
\end{aestheticbox}

\section{The mental picture before the symbols}

Before any Greek letters, fix the picture. We are comparing two functions of $n$,
say your running time $T(n)$ and some clean reference function $g(n)$ like $n^2$ or
$n\log n$. We ask one of three questions:

\begin{itemize}
  \item Does $T$ grow \emph{no faster than} $g$? (An \textbf{upper} bound: $O$.)
  \item Does $T$ grow \emph{no slower than} $g$? (A \textbf{lower} bound: $\Omega$.)
  \item Does $T$ grow \emph{at the same rate} as $g$? (A \textbf{tight} bound: $\Theta$.)
\end{itemize}

Two ideas make these precise, and they recur in every single definition below, so
let me name them now.

\textbf{Idea one: a fudge constant.} Because we refuse to care about constant factors,
we are allowed to multiply $g$ by any positive constant we like to make it fit. So
``$T$ grows no faster than $g$'' really means ``$T(n) \le c\,g(n)$ for some constant
$c$ we get to choose.'' The $c$ absorbs all the machine-dependent grime.

\textbf{Idea two: eventually.} Because small inputs are noise, we only require the
inequality to hold once $n$ is large enough --- past some threshold $n_0$. Below
$n_0$ the functions can cross, wobble, do whatever they like. We care about the tail.

\begin{focusbox}
Every asymptotic statement is of the form: \emph{``there exist constants (a fudge
factor $c$ and a threshold $n_0$) such that for all $n \ge n_0$, some inequality
between $T(n)$ and $g(n)$ holds.''} Memorize that skeleton. The five notations
differ only in \emph{which} inequality and whether the constant is ``some'' or
``every.''
\end{focusbox}

\section{Big-O: the upper bound}

We start with the one you will use most. Informally, $T(n) = O(g(n))$ means $g$ is a
\emph{ceiling} on $T$'s growth: past some point, $T$ never rises above a constant
multiple of $g$.

\begin{formalbox}
\begin{definition}[Big-O]\label{ch03:def-bigo}
Let $g:\N\to\R_{\ge 0}$. We write
\[
  O\bigl(g(n)\bigr) = \set[\big]{\, f(n) : \exists\, c>0,\, n_0>0
  \text{ such that } 0 \le f(n) \le c\,g(n) \text{ for all } n \ge n_0 \,}.
\]
We say ``$f(n)$ is $O(g(n))$'' or write $f(n)=O(g(n))$ to mean $f \in O(g(n))$.
\end{definition}
\end{formalbox}

\begin{trapbox}
Strictly speaking $O(g(n))$ is a \emph{set} of functions, so the cleanest reading
of $f(n)=O(g(n))$ is ``$f$ is a member of the set $O(g(n))$.'' The equals sign here
is a long-standing abuse of notation --- it is \emph{not} symmetric. You may write
$2n^2 = O(n^2)$, but never $O(n^2) = 2n^2$. Read ``$=$'' in asymptotic statements as
the word ``is,'' and only ever from left to right. We follow CLRS's convention and
keep the equals sign, but keep this one-way reading in your head.
\end{trapbox}

Let us draw it. The claim ``$f(n) = O(g(n))$'' says that once we scale $g$ up by the
right constant $c$, the curve $c\,g(n)$ rises above $f(n)$ and stays there.

\begin{center}
\begin{tikzpicture}[scale=1.0]
  \draw[->] (0,0) -- (6.2,0) node[right] {$n$};
  \draw[->] (0,0) -- (0,4.2) node[above] {value};
  % cg(n): the ceiling
  \draw[thick,sourceblue,domain=0.2:6,smooth,variable=\x] plot ({\x},{0.55*\x}) node[right]{$c\,g(n)$};
  % f(n): wobbly, stays under after n0
  \draw[thick,trapred,domain=0.2:6,smooth,variable=\x]
        plot ({\x},{0.32*\x + 0.7*sin(\x r) + 0.6}) node[right]{$f(n)$};
  % n0 marker
  \draw[dashed,gray] (2.3,0) -- (2.3,3.4);
  \node[below] at (2.3,0) {$n_0$};
  \fill[gray!15] (2.3,0) rectangle (6,3.6);
  \node[gray] at (4.3,0.4) {\small for all $n\ge n_0$: $f(n)\le c\,g(n)$};
\end{tikzpicture}
\end{center}

To the left of $n_0$, $f$ is even allowed to poke above $c\,g$. We do not care. To the
right --- the shaded region, the tail that matters --- $f$ is trapped underneath.

\subsection{Proving a Big-O bound: find $c$ and $n_0$}

A Big-O proof is a scavenger hunt for two numbers, $c$ and $n_0$, that make the
definition true. You do not need the \emph{best} constants, just \emph{some} that work.

\begin{example}\label{ch03:ex-poly-bigo}
\emph{Claim:} $\frac12 n^2 + \frac52 n - 3 = O(n^2)$ --- our insertion-sort count is
quadratic.

\emph{Find $c$ and $n_0$.} We want $\frac12 n^2 + \frac52 n - 3 \le c\,n^2$ for all
large $n$. A clean trick: when $n \ge 1$, every term is bounded by a multiple of the
\emph{largest} power, $n^2$, because $n \le n^2$ and $1 \le n^2$. So for $n \ge 1$,
\[
  \tfrac12 n^2 + \tfrac52 n - 3
  \;\le\; \tfrac12 n^2 + \tfrac52 n
  \;\le\; \tfrac12 n^2 + \tfrac52 n^2
  \;=\; 3 n^2 .
\]
(We dropped the $-3$ because subtracting only helps an upper bound, then replaced
$\frac52 n$ by the larger $\frac52 n^2$.) So $c = 3$ and $n_0 = 1$ work. Done.
\hfill$\blacksquare$
\end{example}

Notice the strategy: \emph{keep the dominant term, and bound every smaller term by it.}
That single move proves almost every elementary Big-O bound you will ever need.

\begin{example}\label{ch03:ex-cubic-not}
\emph{Claim:} $n^3 \ne O(n^2)$ --- a cubic is genuinely not quadratic.

\emph{Why no constants can work.} Suppose, for contradiction, that $n^3 \le c\,n^2$
for all $n \ge n_0$. Divide both sides by $n^2$ (positive): $n \le c$ for all
$n \ge n_0$. But $c$ is a fixed number and $n$ marches to infinity, so as soon as
$n > c$ the inequality breaks. No single $c$ can cap $n$. Contradiction. \hfill$\blacksquare$
\end{example}

\begin{walkbox}
Here is one to chew on while walking. Is $2^{n+1} = O(2^n)$? Is $2^{2n} = O(2^n)$?
One of these is true and one is wildly false. Try to write each as a constant times
$2^n$ and see which constant stays put and which one secretly depends on $n$.
(Hint: $2^{n+1} = 2\cdot 2^n$, but $2^{2n} = (2^n)^2$.)
\end{walkbox}

\section{Big-Omega: the lower bound}

$O$ gives a ceiling. \textbf{Big-Omega}, written $\Omega$, gives a \emph{floor}: it
says $T$ grows \emph{at least} as fast as $g$. Flip the inequality in the definition
of $O$ and you have it.

\begin{formalbox}
\begin{definition}[Big-Omega]\label{ch03:def-bigomega}
\[
  \Omega\bigl(g(n)\bigr) = \set[\big]{\, f(n) : \exists\, c>0,\, n_0>0
  \text{ such that } 0 \le c\,g(n) \le f(n) \text{ for all } n \ge n_0 \,}.
\]
\end{definition}
\end{formalbox}

Where $O$ traps $f$ \emph{under} $c\,g$, $\Omega$ traps $f$ \emph{above} $c\,g$.
$O$ is for ``no worse than''; $\Omega$ is for ``no better than.'' A guarantee that an
algorithm is \emph{slow} (a lower bound on running time) is an $\Omega$ statement; a
promise that it is \emph{fast} (an upper bound) is an $O$ statement.

\begin{example}\label{ch03:ex-omega}
\emph{Claim:} $\frac12 n^2 + \frac52 n - 3 = \Omega(n^2)$.

We need $c\,n^2 \le \frac12 n^2 + \frac52 n - 3$ eventually. The $+\frac52 n$ only
helps us (it makes the right side bigger), so it is enough to handle the
$\frac12 n^2 - 3$ part. For $n \ge 4$ we have $\frac14 n^2 \ge \frac14\cdot 16 = 4 \ge 3$,
so $\frac14 n^2$ alone can soak up the $-3$. Thus for $n \ge 4$,
\[
  \tfrac12 n^2 + \tfrac52 n - 3 \;\ge\; \tfrac12 n^2 - 3 \;\ge\; \tfrac12 n^2 - \tfrac14 n^2 = \tfrac14 n^2 .
\]
Thus $c = \frac14$, $n_0 = 4$ work. \hfill$\blacksquare$
\end{example}

The pattern mirrors Big-O: for a \emph{lower} bound, \emph{keep the dominant term and
make sure the negative junk is small enough to be eaten by a fraction of it}.

\begin{trapbox}
$\Omega$ does \emph{not} mean ``the worst case.'' People conflate them constantly.
``Worst case'' picks which \emph{input} to analyze; $\Omega$ and $O$ describe how a
function \emph{grows}. You can put an $O$ bound on the worst case (``never slower
than \dots'') and an $\Omega$ bound on the best case (``never faster than \dots'').
Input choice and bound direction are two independent knobs. Mixing them up produces
sentences that sound rigorous and mean nothing.
\end{trapbox}

\section{Big-Theta: the tight bound}

The crown jewel. $T(n) = \Theta(g(n))$ means $T$ grows \emph{exactly} like $g$ ---
sandwiched between two constant multiples of $g$ from both sides. It is $O$ \emph{and}
$\Omega$ at once.

\begin{formalbox}
\begin{definition}[Big-Theta]\label{ch03:def-bigtheta}
\[
  \Theta\bigl(g(n)\bigr) = \set[\big]{\, f(n) : \exists\, c_1>0,\, c_2>0,\, n_0>0
  \text{ s.t. } 0 \le c_1 g(n) \le f(n) \le c_2 g(n) \text{ for all } n \ge n_0 \,}.
\]
\end{definition}
\end{formalbox}

\begin{center}
\begin{tikzpicture}[scale=1.0]
  \draw[->] (0,0) -- (6.2,0) node[right] {$n$};
  \draw[->] (0,0) -- (0,4.2) node[above] {value};
  \draw[thick,sourceblue,domain=0.2:6,smooth,variable=\x] plot ({\x},{0.58*\x}) node[right]{$c_2 g(n)$};
  \draw[thick,targetgreen,domain=0.2:6,smooth,variable=\x] plot ({\x},{0.22*\x}) node[right]{$c_1 g(n)$};
  \draw[thick,trapred,domain=0.2:6,smooth,variable=\x]
        plot ({\x},{0.38*\x + 0.18*sin(2*\x r)}) node[right]{$f(n)$};
  \draw[dashed,gray] (1.8,0) -- (1.8,3.5);
  \node[below] at (1.8,0) {$n_0$};
  \node[gray] at (4.2,0.35) {\small $f$ is sandwiched for $n\ge n_0$};
\end{tikzpicture}
\end{center}

The picture is the whole story: past $n_0$, the curve $f(n)$ is pinned between the two
straight lines $c_1 g(n)$ and $c_2 g(n)$. It cannot escape upward (that is the $O$ half)
and cannot escape downward (that is the $\Omega$ half). The following theorem makes the
``$\Theta$ = $O$ and $\Omega$'' slogan official.

\begin{formalbox}
\begin{theorem}[Theta as the conjunction of $O$ and $\Omega$]\label{ch03:thm-theta-split}
For any functions $f$ and $g$,
\[
  f(n) = \Theta(g(n)) \quad\text{if and only if}\quad f(n) = O(g(n)) \text{ and } f(n) = \Omega(g(n)).
\]
\end{theorem}
\end{formalbox}

\begin{formalbox}
\emph{Proof.} ($\Rightarrow$) If $f = \Theta(g)$ then there are $c_1,c_2,n_0$ with
$c_1 g(n) \le f(n) \le c_2 g(n)$ for $n\ge n_0$. The right inequality alone, with
$c = c_2$, witnesses $f = O(g)$; the left inequality alone, with $c = c_1$, witnesses
$f = \Omega(g)$.

($\Leftarrow$) Suppose $f = O(g)$ via $(c_2, n_2)$ and $f = \Omega(g)$ via $(c_1, n_1)$.
Take $n_0 = \max(n_1, n_2)$. For every $n \ge n_0$ \emph{both} witnessing inequalities
hold simultaneously, so $c_1 g(n) \le f(n) \le c_2 g(n)$, which is exactly the
$\Theta$ condition with constants $c_1, c_2$ and threshold $n_0$. \hfill$\blacksquare$
\end{formalbox}

Combining Examples~\ref{ch03:ex-poly-bigo} and~\ref{ch03:ex-omega} through this theorem,
we conclude $\frac12 n^2 + \frac52 n - 3 = \Theta(n^2)$. Insertion sort's worst case is
\emph{tightly} quadratic --- not just bounded above by $n^2$, but genuinely of that order.

\subsection{A theorem you will reuse: polynomials are Theta of their top term}

The two examples generalize into a workhorse fact. It justifies the move you have
already been itching to make: just read off the highest power.

\begin{formalbox}
\begin{theorem}[Degree dominates]\label{ch03:thm-poly}
Let $p(n) = a_d n^d + a_{d-1} n^{d-1} + \cdots + a_1 n + a_0$ be a polynomial of degree
$d$ with leading coefficient $a_d > 0$. Then $p(n) = \Theta(n^d)$.
\end{theorem}
\end{formalbox}

\begin{formalbox}
\emph{Proof sketch.} \textbf{Upper bound.} For $n \ge 1$ each $n^k \le n^d$, so
\[
  p(n) \le \bigl(\,\abs{a_d} + \abs{a_{d-1}} + \cdots + \abs{a_0}\,\bigr) n^d =: c_2\, n^d,
\]
giving $p = O(n^d)$. \textbf{Lower bound.} Factor out $n^d$:
$p(n) = n^d\bigl(a_d + a_{d-1}/n + \cdots + a_0/n^d\bigr)$. The parenthesized factor
tends to $a_d > 0$ as $n\to\infty$, so for all $n$ past some $n_0$ it exceeds
$a_d/2 > 0$. Hence $p(n) \ge (a_d/2)\,n^d$, giving $p = \Omega(n^d)$. By
Theorem~\ref{ch03:thm-theta-split}, $p = \Theta(n^d)$. \hfill$\blacksquare$
\end{formalbox}

\begin{focusbox}
\textbf{Two rules you will use a hundred times.}
\emph{Drop low-order terms} (they are out-grown by the leading term) and
\emph{drop constant factors} (they are absorbed by $c$). So
$7n^3 - 100n^2 + 4n + 9999$ is just $\Theta(n^3)$, and
$\frac{n\log n}{2} + 30n$ is just $\Theta(n\log n)$. The notation exists precisely so
you may make these two moves without apology.
\end{focusbox}

\section{The little ones: little-o and little-omega}

$O$ and $\Omega$ are loose --- they allow $f$ and $g$ to grow at the \emph{same} rate.
Indeed $n^2 = O(n^2)$. The \emph{little} versions are strict: they insist that one
function grows \emph{strictly slower} (little-o) or \emph{strictly faster}
(little-omega) than the other, with room to spare.

The cleanest way to feel the difference: in $O$, \emph{some} constant $c$ works. In
little-o, the bound must hold for \emph{every} positive constant, no matter how tiny.

\begin{formalbox}
\begin{definition}[little-o]\label{ch03:def-littleo}
$f(n) = o(g(n))$ if for \emph{every} constant $c>0$ there exists $n_0>0$ such that
$0 \le f(n) < c\,g(n)$ for all $n \ge n_0$.
\end{definition}
\begin{definition}[little-omega]\label{ch03:def-littleomega}
$f(n) = \omega(g(n))$ if for \emph{every} constant $c>0$ there exists $n_0>0$ such that
$0 \le c\,g(n) < f(n)$ for all $n \ge n_0$.
\end{definition}
\end{formalbox}

That ``for every $c$'' is the whole game. Take $f(n) = 2n$ and $g(n) = n^2$. Is
$2n = o(n^2)$? Pick \emph{any} $c$, say $c = 0.001$. We need $2n < 0.001\,n^2$, i.e.
$2 < 0.001\,n$, i.e. $n > 2000$. So $n_0 = 2001$ works for $c=0.001$. Try $c = 10^{-9}$
and you would just need a bigger $n_0$ --- but one always exists. The smaller the $c$,
the longer you wait, but you always eventually win. That is little-o.

By contrast $2n^2 = O(n^2)$ but $2n^2 \ne o(n^2)$: choose $c = 1$ and you would need
$2n^2 < n^2$, i.e. $2 < 1$, which never holds. A function is never little-o of itself.

\begin{intuitionbox}
A clean way to remember the family, when the limit exists, is to look at the ratio
$f(n)/g(n)$ as $n\to\infty$:
\[
  \lim_{n\to\infty}\frac{f(n)}{g(n)} =
  \begin{cases}
    0 & \Rightarrow\; f = o(g)\quad(\text{$f$ vanishes next to $g$}),\\[2pt]
    \text{a positive constant} & \Rightarrow\; f = \Theta(g)\quad(\text{same rate}),\\[2pt]
    \infty & \Rightarrow\; f = \omega(g)\quad(\text{$f$ runs away from $g$}).
  \end{cases}
\]
Big-$O$ is the union of the first two cases ($0$ or a constant); big-$\Omega$ is the
union of the last two. The limit test is not the official definition (the limit need
not exist), but when it exists it is the fastest way to classify a pair of functions.
\end{intuitionbox}

\begin{example}\label{ch03:ex-logvspoly}
\emph{$\log n = o(n)$.} The ratio $(\log_2 n)/n \to 0$: logarithms crawl while $n$
strides. So any logarithm is little-o of $n$, hence of any positive power $n^\eps$.
Logs are, asymptotically, almost free.
\end{example}

\begin{walkbox}
A genuinely surprising one to ponder on a walk: is $n^{100} = o(2^n)$? It feels like
$n^{100}$ should win --- the exponent is huge. But $2^n$ is exponential and
\emph{every} polynomial is little-o of \emph{every} exponential with base $> 1$.
A clean way to see it: take $\log_2$ of the ratio, $\log_2\!\bigl(n^{100}/2^n\bigr)
= 100\log_2 n - n$. The $-n$ term grows steadily while $100\log_2 n$ crawls, so the
whole expression marches to $-\infty$; hence the ratio itself marches to $0$. Even with
the degree cranked to $100$, the crossover merely waits until around $n = 1000$ and then
$2^n$ pulls ahead and never trails again. Exponentials always win, no matter how large
the polynomial's degree.
\end{walkbox}

\section{An analogy: comparing functions is comparing numbers}

The five notations line up neatly against the relations you already know between plain
numbers:

\begin{center}
\begin{tabular}{ccl}
\toprule
Functions & Numbers & Reading \\
\midrule
$f = O(g)$       & $a \le b$ & $f$ grows no faster than $g$ \\
$f = \Omega(g)$  & $a \ge b$ & $f$ grows no slower than $g$ \\
$f = \Theta(g)$  & $a = b$   & $f$ grows at the same rate as $g$ \\
$f = o(g)$       & $a < b$   & $f$ grows strictly slower than $g$ \\
$f = \omega(g)$  & $a > b$   & $f$ grows strictly faster than $g$ \\
\bottomrule
\end{tabular}
\end{center}

This analogy is more than a memory aid; many properties carry straight over.
\emph{Transitivity} holds: if $f = O(g)$ and $g = O(h)$ then $f = O(h)$, just as
$a\le b$ and $b\le c$ give $a \le c$. \emph{Reflexivity} holds for the big ones:
$f = O(f)$, $f=\Theta(f)$, $f=\Omega(f)$ (but \emph{not} for the little ones --- nothing
is $o$ of itself, just as no number is $<$ itself). \emph{Symmetry} holds only for
$\Theta$: $f=\Theta(g)$ iff $g=\Theta(f)$, mirroring $a=b$ iff $b=a$. And the little/big
pair are \emph{transposes}: $f = o(g)$ iff $g = \omega(f)$, exactly like $a<b$ iff $b>a$.

\begin{trapbox}
The analogy has \emph{one crucial leak}: numbers obey \emph{trichotomy} --- any two
numbers satisfy exactly one of $a<b$, $a=b$, $a>b$. Functions do not. Two functions can
be \emph{incomparable}: neither $f = O(g)$ nor $f = \Omega(g)$. For instance,
$f(n) = n$ and $g(n) = n^{1 + \sin n}$ keep overtaking each other forever, so neither
bounds the other. In practice the functions you meet (polynomials, logs, exponentials)
are all comparable, so the trichotomy intuition usually serves --- but know that it can
break.
\end{trapbox}

\section{A field guide to common growth rates}

Here are the functions you will meet over and over, sorted slowest-growing (best) to
fastest-growing (worst). Each line is little-o of the line below it.

\begin{center}
\begin{tabular}{lll}
\toprule
Class & Name & Where it shows up \\
\midrule
$\Theta(1)$        & constant      & array index, hash lookup (ideal) \\
$\Theta(\log n)$   & logarithmic   & binary search, balanced-tree height \\
$\Theta(n)$        & linear        & one scan of the input \\
$\Theta(n\log n)$  & linearithmic  & mergesort, heapsort, good comparison sorts \\
$\Theta(n^2)$      & quadratic     & insertion/selection sort, all-pairs naively \\
$\Theta(n^3)$      & cubic         & naive matrix multiply, Floyd--Warshall \\
$\Theta(2^n)$      & exponential   & brute-force subset search \\
$\Theta(n!)$       & factorial     & brute-force permutation search (TSP) \\
\bottomrule
\end{tabular}
\end{center}

Figure~\ref{ch03:fig-growth} draws the most common classes on one set of axes. Notice how
brutally the ordering separates them once $n$ leaves the cramped corner near the origin:
$\log n$ flattens almost to the floor, $n$ and $n\log n$ rise gently, $n^2$ curls upward,
and $2^n$ shoots off the top of the page almost immediately.

\begin{figure}[H]
\centering
\begin{tikzpicture}[scale=1.0]
  \draw[->] (0,0) -- (6.4,0) node[right] {$n$};
  \draw[->] (0,0) -- (0,4.6) node[above] {steps};
  % logarithmic: very flat
  \draw[thick,targetgreen,domain=0.4:6,smooth,variable=\x]
        plot ({\x},{0.55*ln(\x+1)}) node[right]{$\log n$};
  % linear
  \draw[thick,sourceblue,domain=0:5.6,smooth,variable=\x]
        plot ({\x},{0.42*\x}) node[right]{$n$};
  % n log n
  \draw[thick,quotientorange,domain=0.4:4.6,smooth,variable=\x]
        plot ({\x},{0.30*\x*ln(\x+1)}) node[right]{$n\log n$};
  % quadratic
  \draw[thick,structurepurple,domain=0:2.95,smooth,variable=\x]
        plot ({\x},{0.46*\x*\x}) node[above right]{$n^2$};
  % exponential: shoots up
  \draw[thick,trapred,domain=0:2.45,smooth,variable=\x]
        plot ({\x},{0.18*exp(1.15*\x)}) node[right]{$2^n$};
\end{tikzpicture}
\caption{Common growth rates on shared axes. The vertical order at the right edge is the
asymptotic order: slower-growing curves stay low, faster-growing curves climb away. Each
curve is eventually little-o of every curve above it.}
\label{ch03:fig-growth}
\end{figure}

The chasm between the polynomial classes (everything down to $n^3$) and the exponential
ones is the single most important divide in all of algorithms; an entire late chapter
on NP-completeness is about which problems we can keep on the polynomial side. To feel
the chasm, put real seconds on it. Suppose your machine does $10^9$ basic steps per
second:

\begin{center}
\begin{tabular}{rrrr}
\toprule
$n$ & $n\log_2 n$ & $n^2$ & $2^n$ \\
\midrule
$10$  & $33$ ns       & $100$ ns    & $1024$ ns \\
$20$  & $86$ ns       & $400$ ns    & $1.0$ ms \\
$50$  & $282$ ns      & $2.5$ \textmu s & $13$ days \\
$100$ & $664$ ns      & $10$ \textmu s  & $4\times10^{13}$ yr \\
$10^6$ & $20$ ms      & $17$ min    & --- (heat death) \\
\bottomrule
\end{tabular}
\end{center}

At $n = 100$, $n\log n$ finishes before you blink and $2^n$ outlasts the universe by
many orders of magnitude. \emph{This} is why growth rate, not the constant, is the thing
to fight over.

\begin{aestheticbox}
Logarithms deserve special affection. Because $\log n$ grows so slowly, an
$O(\log n)$ operation is nearly constant for any input you will ever see: $\log_2$ of a
billion is about $30$, and $\log_2$ of the number of atoms in the observable universe is
under $270$. When you turn an $O(n)$ search into an $O(\log n)$ one --- the leap from
scanning a list to binary-searching a sorted array, or to walking a balanced tree --- you
are not shaving a little time. You are buying the right to make your data a billion times
bigger for the cost of a handful of extra steps. Whole data structures in later chapters
exist to manufacture that logarithm.
\end{aestheticbox}

\begin{historybox}
Big-$O$ notation was born in number theory, not computer science: Paul Bachmann
introduced the symbol in $1894$ and Edmund Landau popularized it (the symbols are still
sometimes called ``Landau notation''). Donald Knuth, in the $1970$s, championed its use
for algorithm analysis and, finding the literature sloppy, pinned down the matching
$\Omega$ and $\Theta$ with the precise meanings we use today. This chapter mirrors
\textbf{CLRS Chapter 3, ``Characterizing Running Times''} (Section 3.1 on $O$, $\Omega$,
$\Theta$ and Section 3.2 on standard notations and common functions). When you open CLRS,
the set-based definitions and the asymmetric ``$=$'' will already be old friends; CLRS
adds the wrinkle of asymptotic notation \emph{inside} equations (e.g. $2n^2 + \Theta(n)$),
which is just these same ideas used as a noun mid-formula.
\end{historybox}

\section{Manipulating asymptotics: the algebra of growth}

You rarely build a bound from $c$ and $n_0$ in the wild. Instead you compose small
facts with a few algebraic rules. Here are the ones that do the daily work; each follows
directly from the definitions, and proving one or two yourself is the best exercise of
the chapter.

\begin{formalbox}
\begin{lemma}[Working rules for $O$]\label{ch03:lem-rules}
Let $f_1 = O(g_1)$ and $f_2 = O(g_2)$. Then:
\begin{enumerate}
  \item \textbf{Sum / max:} $f_1 + f_2 = O\bigl(\max(g_1,g_2)\bigr) = O(g_1 + g_2)$.
  \item \textbf{Product:} $f_1 \cdot f_2 = O(g_1 \cdot g_2)$.
  \item \textbf{Scaling:} for any constant $a>0$, $a\cdot f_1 = O(g_1)$.
  \item \textbf{Transitivity:} if also $g_1 = O(h)$, then $f_1 = O(h)$.
\end{enumerate}
The same four rules hold verbatim with $O$ replaced by $\Omega$ or by $\Theta$.
\end{lemma}
\end{formalbox}

Rule 1 is the one beginners trip over, so let us prove it and see why ``the bigger one
wins.''

\begin{formalbox}
\emph{Proof of the sum rule.} We are given $f_1(n) \le c_1 g_1(n)$ for $n\ge n_1$ and
$f_2(n) \le c_2 g_2(n)$ for $n \ge n_2$. Let $n_0 = \max(n_1,n_2)$ and $c = c_1 + c_2$.
For $n \ge n_0$ both bounds hold, so
\[
  f_1(n) + f_2(n) \;\le\; c_1 g_1(n) + c_2 g_2(n)
  \;\le\; (c_1 + c_2)\max\bigl(g_1(n), g_2(n)\bigr) = c\cdot \max(g_1,g_2)(n).
\]
Thus $f_1 + f_2 = O(\max(g_1,g_2))$. \hfill$\blacksquare$
\end{formalbox}

\begin{focusbox}
\textbf{The sum rule is why you keep only the dominant term.} A program that does an
$O(n)$ phase, then an $O(n^2)$ phase, then an $O(n\log n)$ phase runs in
$O(n + n^2 + n\log n) = O(n^2)$ --- the slowest phase swallows the rest. When you read a
running time, your eye should jump straight to the heaviest term and ignore the others.
\end{focusbox}

\begin{example}\label{ch03:ex-nested-loops}
A doubly nested loop where the inner loop's length grows with the outer index:
\[
  \sum_{i=1}^{n}\; \sum_{j=1}^{i} 1 \;=\; \sum_{i=1}^{n} i \;=\; \frac{n(n+1)}{2}
  \;=\; \Theta(n^2).
\]
This is the cost shape of insertion sort and of many ``compare each pair'' routines:
the inner work is $1 + 2 + \cdots + n$, a triangle of area $\approx n^2/2$, which the
notation cleanly reports as $\Theta(n^2)$. (The appendix on summations derives
$\sum i = n(n+1)/2$ from scratch if you want the gory details.)

Figure~\ref{ch03:fig-triangle} makes the ``triangle'' literal: stack one dot for the
inner work at $i=1$, two dots at $i=2$, and so on. The dots fill half of an
$n\times n$ square, so their count is about $n^2/2$ --- a constant ($\tfrac12$) times
$n^2$, which is exactly what $\Theta(n^2)$ asserts.

\begin{figure}[H]
\centering
\begin{tikzpicture}[scale=0.42]
  \def\m{8}
  % full n x n square outline
  \draw[gray!45] (0,0) rectangle (\m,\m);
  % the triangle of dots: row r (from bottom) has r dots
  \foreach \r in {1,...,\m}{
    \foreach \col in {1,...,\r}{
      \fill[sourceblue] (\col-0.5, \r-0.5) circle (0.18);
    }
  }
  % diagonal showing the half
  \draw[trapred,thick] (0,0) -- (\m,\m);
  \node[gray] at (\m+2.6,\m-0.6) {\small full square: $n^2$};
  \node[sourceblue] at (\m+2.6,2.2) {\small dots: $\tfrac{n(n+1)}{2}$};
  \node[below=2pt] at (\m/2,0) {\small outer index $i = 1,2,\ldots,n$};
\end{tikzpicture}
\caption{The work of nested loops as a triangle. Column $i$ holds $i$ units of inner
work; the dots fill just over half the $n\times n$ square, so the total is
$\tfrac{n(n+1)}{2} = \Theta(n^2)$.}
\label{ch03:fig-triangle}
\end{figure}
\end{example}

\begin{labbox}
\textbf{Trace the growth yourself.} Take the function $T(n) = 3n^2 + 50n + 200$.
\begin{enumerate}
  \item Build a table of $T(n)$, $3n^2$, and the ratio $T(n)/n^2$ for
        $n = 1, 10, 100, 1000$. Watch the ratio slide toward $3$.
  \item That limiting ratio of $3$ is your $c_1$/$c_2$ neighborhood: find a single
        $c_2$ with $T(n) \le c_2 n^2$ for all $n\ge 1$ (try $c_2 = 253$ bluntly, then
        sharpen it). Then find $c_1$ and $n_0$ with $c_1 n^2 \le T(n)$ (try $c_1 = 3$).
  \item Conclude $T(n) = \Theta(n^2)$ by exhibiting your $c_1, c_2, n_0$. Notice how the
        constants $50$ and $200$ never appear in the answer --- they only pushed $n_0$
        a little to the right.
\end{enumerate}
This is the entire skill: turn a messy formula into one Greek symbol by finding
witnesses.
\end{labbox}

\section{Putting it to work: bounding an algorithm directly}

Let us close the loop by reading a running time straight off pseudocode, the way you
will for the rest of the book. Here is a routine that, given an array $A$ of length $n$,
checks whether it contains any duplicated value by brute force.

\begin{algorithm}[H]
\caption{\textsc{HasDuplicate}$(A, n)$}
\KwIn{An array $A[1\ldots n]$ of $n$ values.}
\KwOut{\textsc{true} if some value appears twice, else \textsc{false}.}
\For{$i \leftarrow 1$ \KwTo $n-1$}{
  \For{$j \leftarrow i+1$ \KwTo $n$}{
    \If{$A[i] = A[j]$}{
      \Return \textsc{true}\;
    }
  }
}
\Return \textsc{false}\;
\end{algorithm}

\textbf{Counting the work.} The body of the inner loop (the comparison) is the basic
operation. For a fixed $i$, the inner loop runs $j$ from $i+1$ to $n$, that is $n - i$
times. Summing over all outer iterations,
\[
  \sum_{i=1}^{n-1} (n - i) \;=\; (n-1) + (n-2) + \cdots + 1 \;=\; \frac{(n-1)n}{2}.
\]
By Theorem~\ref{ch03:thm-poly} this is $\Theta(n^2)$ --- \emph{in the worst case}, when
no early \textsc{return} fires (no duplicate, or the only duplicate is the very last
pair). So the worst-case running time is $\Theta(n^2)$.

\textbf{What about best case?} If $A[1] = A[2]$, the routine returns after a single
comparison: best-case cost is $\Theta(1)$. So we can say the running time is $O(n^2)$
(it is never worse than quadratic) and $\Omega(1)$ (it always does at least constant
work) --- but we \emph{cannot} wrap the whole thing in a single $\Theta$, because best
and worst cases have different growth rates. Many algorithms sit in exactly this
position, and the notation lets us say as much as is true and no more.

\begin{trapbox}
``$\Theta$ of the worst case'' and ``$\Theta$ of the running time'' are different
claims. \textsc{HasDuplicate}'s \emph{worst case} is $\Theta(n^2)$ --- a tight statement
about one specific input family. Its \emph{running time} overall is only $O(n^2)$ and
$\Omega(1)$, because the cost genuinely varies with the input. Reserve a single $\Theta$
for when best and worst cases share a growth rate (as in mergesort, always
$\Theta(n\log n)$). When they differ, say $O$ for the ceiling and $\Omega$ for the floor,
and resist the urge to collapse them.
\end{trapbox}

\begin{complexitybox}
\textbf{The five notations at a glance.}
\begin{center}
\begin{tabular}{lll}
\toprule
Notation & Meaning & Number analogue \\
\midrule
$f = O(g)$       & $\le c\,g$ eventually, some $c$        & $\le$ \\
$f = \Omega(g)$  & $\ge c\,g$ eventually, some $c$        & $\ge$ \\
$f = \Theta(g)$  & $c_1 g \le f \le c_2 g$ eventually     & $=$ \\
$f = o(g)$       & $< c\,g$ eventually, \emph{every} $c$  & $<$ \\
$f = \omega(g)$  & $> c\,g$ eventually, \emph{every} $c$  & $>$ \\
\bottomrule
\end{tabular}
\end{center}
\textbf{Every} bound is ``$\exists$ constants, $\forall\, n \ge n_0$, inequality.''
For polynomials: $\Theta$ of the top power. For sums: keep the dominant term. For
products: multiply the bounds. Logs $\prec$ powers $\prec$ exponentials $\prec$
factorials, each gap a little-o.
\end{complexitybox}

\section{Exercises}

\begin{exercise}\label{ch03:ex-warm-classify}
For each pair, state whether $f = O(g)$, $f = \Omega(g)$, both (i.e. $\Theta$), and
whether either little version holds:
(a) $f = 5n$, $g = n^2$;\;
(b) $f = n^2 + n$, $g = 100n^2$;\;
(c) $f = 2^n$, $g = 3^n$;\;
(d) $f = \log_2 n$, $g = \log_{10} n$;\;
(e) $f = n$, $g = n\log n$.
\end{exercise}

\begin{exercise}\label{ch03:ex-find-constants}
Prove that $3n^2 + 10n + 7 = O(n^2)$ by exhibiting explicit constants $c$ and $n_0$ and
verifying the inequality for all $n \ge n_0$. (Any valid pair earns full marks; aim for
small, clean numbers.)
\end{exercise}

\begin{exercise}\label{ch03:ex-omega-proof}
Prove that $\frac{n^2}{2} - 3n = \Theta(n^2)$. You must supply \emph{both} an upper
bound ($c_2, n_0$) and a lower bound ($c_1, n_0$); take care that the $-3n$ does not
sink the lower bound.
\end{exercise}

\begin{exercise}\label{ch03:ex-disprove}
Show that $n^2 \ne O(n)$ by the contradiction method of
Example~\ref{ch03:ex-cubic-not}: assume $n^2 \le c n$ for $n \ge n_0$ and derive an
impossibility.
\end{exercise}

\begin{exercise}\label{ch03:ex-transitivity}
Using only the definition of $O$ (constants $c$ and $n_0$), prove the transitivity rule:
if $f = O(g)$ and $g = O(h)$, then $f = O(h)$. Be explicit about how you combine the two
sets of witnessing constants.
\end{exercise}

\begin{exercise}\label{ch03:ex-symmetry}
Prove that $f(n) = \Theta(g(n))$ if and only if $g(n) = \Theta(f(n))$ (the symmetry of
$\Theta$). Then give a concrete $f, g$ showing that $f = O(g)$ does \emph{not} imply
$g = O(f)$, so $O$ is \emph{not} symmetric.
\end{exercise}

\begin{exercise}\label{ch03:ex-littleo-limit}
Use the limit characterization to decide each: (a) is $n\log n = o(n^2)$? (b) is
$\sqrt{n} = o(n)$? (c) is $n = o(n + \sqrt n)$? For each, compute
$\lim_{n\to\infty} f(n)/g(n)$ and read off the answer.
\end{exercise}

\begin{exercise}\label{ch03:ex-nested}
A triple-nested loop runs $i$ from $1$ to $n$, $j$ from $1$ to $i$, and $k$ from $1$ to
$j$, doing constant work in the innermost body. Write the total work as a double sum,
evaluate it, and give the tight $\Theta$ bound. (You may use $\sum_{i=1}^n i^2 =
\frac{n(n+1)(2n+1)}{6}$ from the summations appendix.)
\end{exercise}

\begin{exercise}\label{ch03:ex-exp-beats-poly}
\emph{(Challenge.)} Prove that $n^k = o(2^n)$ for every fixed positive integer $k$.
\emph{Hint:} consider $\log_2$ of the ratio, $\log_2\!\bigl(n^k / 2^n\bigr) = k\log_2 n - n$,
and argue that this tends to $-\infty$, so the ratio tends to $0$.
\end{exercise}

\begin{exercise}\label{ch03:ex-incomparable}
\emph{(Challenge.)} Give two functions $f, g : \N \to \R_{\ge 0}$ such that
\emph{neither} $f = O(g)$ \emph{nor} $g = O(f)$ holds, and explain in one or two
sentences why each $O$ statement fails. (This shows asymptotic comparison is not a total
order --- trichotomy can fail.)
\end{exercise}

\begin{exercise}\label{ch03:ex-merge-vs-insertion}
\emph{(Synthesis.)} Insertion sort runs in $\Theta(n^2)$ worst case; mergesort runs in
$\Theta(n\log n)$ always. Suppose on your machine insertion sort costs exactly $2n^2$
nanoseconds and mergesort costs exactly $50 n\log_2 n$ nanoseconds. Find, by hand or by a
small table, the approximate input size $n^\ast$ beyond which mergesort wins, and explain
why the existence of such a crossover is guaranteed by the asymptotics no matter how the
two constants ($2$ and $50$) had been chosen.
\end{exercise}
